\documentclass[../../main.tex]{subfiles}

\begin{document}

\chapter{Operasi Aritmatika}
\label{chap:aritmatika}

% =========================================================
% OUTCOME PEMBELAJARAN KHUSUS BAB
% =========================================================

\begin{subcpmkbox}
\vspace{1.1em}
\SubCPMK{3.1} --- Menulis program dengan instruksi ADD, ADC, INC,
SUB, SBB, DEC, MUL, dan DIV serta menjelaskan pengaruhnya terhadap
\eng{flag}.

\vspace{0.4em}
\textbf{CPMK yang didukung:} \CPMK{03} \quad\textbullet\quad
\textbf{CPL yang didukung:} \CPL{03}
\end{subcpmkbox}

\begin{tujuanbox}
\vspace{1.1em}
Setelah mempelajari bab ini, mahasiswa mampu:
\begin{enumerate}
    \item menggunakan \asm{ADD}, \asm{ADC}, \asm{INC}, \asm{SUB},
          \asm{SBB}, dan \asm{DEC} pada operand 8 bit dan 16 bit;
    \item menjelaskan pengaruh operasi aritmatika terhadap
          \flag{CF}, \flag{OF}, \flag{ZF}, \flag{SF}, \flag{AF}, dan
          \flag{PF};
    \item membedakan makna \flag{CF} untuk bilangan tidak bertanda
          dan \flag{OF} untuk bilangan bertanda;
    \item menyusun penjumlahan dan pengurangan lebih dari 16 bit
          dengan \asm{ADC} dan \asm{SBB};
    \item menggunakan \asm{MUL} dan \asm{DIV} sesuai aturan operand
          implisit;
    \item memeriksa dan menghindari \eng{divide overflow};
    \item memverifikasi hasil operasi dengan \eng{debugger}.
\end{enumerate}
\end{tujuanbox}

\begin{materibox}
\vspace{1.1em}
\begin{enumerate}
    \item Penjumlahan: \asm{ADD}, \asm{ADC}, dan \asm{INC}.
    \item Pengurangan: \asm{SUB}, \asm{SBB}, dan \asm{DEC}.
    \item Makna \eng{flag} hasil aritmatika.
    \item Aritmatika bertingkat dengan \asm{ADC} dan \asm{SBB}.
    \item Perkalian \asm{MUL} dan aturan operandnya.
    \item Pembagian \asm{DIV} dan \eng{divide overflow}.
    \item Verifikasi hasil dengan DEBUG dan emu8086.
\end{enumerate}
\end{materibox}

\subfile{section-01}
\subfile{section-02}
\subfile{section-03}

\section{Contoh}
\label{sec:contoh-bab10}

Dua contoh berikut memperlihatkan operasi aritmatika lengkap dan
pemeriksaan hasilnya. Contoh pertama adalah program utuh yang
menyimpan hasil ke segmen data. Contoh kedua memperlihatkan
perkalian dan pembagian 16 bit dengan pasangan \reg{DX}:\reg{AX}.

\subsection{Contoh 1: Operasi Aritmatika Lengkap}

Program pada Kode~\ref{lst:operasi-aritmatika} memakai \asm{ADD},
\asm{ADC}, \asm{SUB}, \asm{SBB}, \asm{INC}, \asm{DEC}, \asm{MUL},
dan \asm{DIV} dalam satu program.

\lstinputlisting[
    language=MASM,
    caption={Program \file{operasi-aritmatika.asm} untuk mengamati
    operasi aritmatika dan \eng{flag}.},
    label={lst:operasi-aritmatika}
]{\CodeDir/operasi-aritmatika.asm}

Program tidak mencetak keluaran. Hasil setiap operasi disimpan pada
segmen data dan diperiksa dengan DEBUG atau emu8086. Bagian penting
program tersebut adalah sebagai berikut.

\begin{itemize}
    \item \lstasm|ADD AX, b_lo| menjumlahkan word rendah. Karena
          hasilnya melewati 16 bit, \flag{CF} bernilai 1.
    \item \lstasm|ADC AX, b_hi| menjumlahkan word tinggi ditambah
          \flag{CF}, sehingga hasil 32 bit menjadi
          \hex{0004}:\hex{0000}.
    \item \lstasm|SUB AX, d_lo| menghasilkan pinjaman, lalu
          \lstasm|SBB AX, d_hi| melanjutkan pengurangan dengan
          pinjaman tersebut.
    \item \lstasm|MUL BL| mengalikan \reg{AL} dengan \reg{BL} dan
          menaruh hasil 16 bit pada \reg{AX}.
    \item \lstasm|CMP AH, BL| memeriksa apakah hasil bagi \asm{DIV}
          akan muat pada \reg{AL}. Bila tidak, program melompat ke
          label \code{luapan}.
\end{itemize}

Hasil yang diharapkan dirangkum pada Tabel~\ref{tab:hasil-program}.

\begin{table}[H]
    \centering
    \caption{Hasil operasi pada \file{operasi-aritmatika.asm}.}
    \label{tab:hasil-program}
    \small
    \begin{tabular}{@{}p{4.4cm}p{2.6cm}p{6.2cm}@{}}
        \toprule
        \textbf{Operasi} & \textbf{Hasil} &
        \textbf{Letak hasil} \\
        \midrule
        \hex{0001FFFF} + \hex{00020001} &
        \hex{0004}:\hex{0000} &
        \code{hasil\_hi}:\code{hasil\_lo} \\
        \addlinespace
        \hex{00050000} $-$ \hex{00020001} &
        \hex{0002}:\hex{FFFF} &
        \code{kurang\_hi}:\code{kurang\_lo} \\
        \addlinespace
        $25 \times 4$ & \hex{0064} & \code{kali} \\
        \addlinespace
        $1000 \div 7$ & \hex{008E}, sisa \hex{0006} &
        \code{bagi}, \code{sisa} \\
        \bottomrule
    \end{tabular}
\end{table}

\subsection{Contoh 2: MUL dan DIV 16 Bit}

Perkalian 16 bit menghasilkan 32 bit pada \reg{DX}:\reg{AX}.
Pembagian 16 bit membagi \reg{DX}:\reg{AX} dengan operand dan
menyimpan hasil bagi di \reg{AX} serta sisa di \reg{DX}.

\begin{lstlisting}[language=MASM, numbers=none]
; MUL 16 bit: 1000 x 20 = 20000 (4E20h)
MOV AX, 1000
MOV BX, 20
MUL BX              ; DX:AX = 0000:4E20h

; DIV 16 bit: 100000 / 1000 = 100
MOV DX, 1
MOV AX, 86A0h       ; DX:AX = 0001:86A0h = 100000
MOV BX, 1000
CMP DX, BX          ; periksa agar hasil bagi muat di AX
JAE luapan
DIV BX              ; AX = 100, DX = 0
luapan:
\end{lstlisting}

Pada pembagian 16 bit, syarat aman adalah \reg{DX} lebih kecil
daripada pembagi. Bila \reg{DX} sama atau lebih besar, hasil bagi
tidak muat pada \reg{AX} dan terjadi \eng{divide overflow}.

\begin{infobox}
\vspace{1.1em}
Nilai \hex{4E20} adalah 20.000 desimal. Angka ini masih muat pada
\reg{AX}, sehingga \reg{DX} bernilai nol. Bila hasil perkalian
melebihi 65.535, \reg{DX} tidak nol dan \flag{CF} = \flag{OF} = 1.
\end{infobox}

\section{Praktikum}
\label{sec:praktikum-bab10}

\begin{praktikumbox}
\vspace{1.1em}
\textbf{Judul:} Menelusuri operasi aritmatika dan memeriksa
hasilnya.\\[0.3em]
\textbf{Sub-CPMK:} 3.1 \quad\textbullet\quad
\textbf{Alokasi waktu:} 100 menit \quad\textbullet\quad
\textbf{Bentuk:} individu
\end{praktikumbox}

\subsection{Tujuan Praktikum}

Setelah menyelesaikan praktikum ini, mahasiswa mampu menjalankan
instruksi aritmatika, mencatat perubahan register dan \eng{flag},
serta memverifikasi hasil operasi dengan DEBUG dan emu8086.

\subsection{Alat dan Bahan}

\begin{itemize}
    \item Komputer dengan emu8086 terpasang.
    \item DOSBox dengan TASM, TLINK, dan DEBUG.
    \item Berkas \file{operasi-aritmatika.asm} dari \file{code/asm/}.
    \item Lembar pengamatan register dan \eng{flag}.
\end{itemize}

\subsection{Langkah Kerja}

\begin{enumerate}
    \item \textbf{Merakit dan menjalankan program.}
    \begin{enumerate}
        \item Buka DOSBox, lalu masuk ke direktori kerja.
        \item Rakit dengan \file{tasm operasi-aritmatika.asm}.
        \item Link dengan \file{tlink operasi-aritmatika.obj}.
        \item Jalankan program. Program tidak mencetak keluaran.
    \end{enumerate}

    \item \textbf{Memeriksa hasil di segmen data.}
    \begin{enumerate}
        \item Jalankan \file{debug operasi-aritmatika.exe}.
        \item Ketik \file{d ds:0000 L 20} untuk melihat isi segmen
              data.
        \item Catat nilai \code{hasil\_lo}, \code{hasil\_hi},
              \code{kurang\_lo}, \code{kurang\_hi}, \code{kali},
              \code{bagi}, dan \code{sisa}.
        \item Bandingkan dengan Tabel~\ref{tab:hasil-program}.
    \end{enumerate}

    \item \textbf{Menelusuri dengan emu8086.}
    \begin{enumerate}
        \item Buka program pada emu8086.
        \item Jalankan \file{Single Step} sampai instruksi
              \lstasm|ADD AX, b_lo|.
        \item Catat \reg{AX}, \reg{BX}, dan \eng{flag} sebelum dan
              sesudah setiap instruksi.
        \item Ulangi untuk \lstasm|ADC AX, b_hi|,
              \lstasm|SUB AX, d_lo|, dan \lstasm|SBB AX, d_hi|.
    \end{enumerate}

    \item \textbf{Menguji MUL dan DIV.}
    \begin{enumerate}
        \item Catat \reg{AX} setelah \lstasm|MUL BL|.
        \item Catat \reg{AL} dan \reg{AH} setelah \lstasm|DIV BL|.
        \item Jelaskan perbedaan letak hasil bagi dan sisa.
    \end{enumerate}

    \item \textbf{Mengubah nilai untuk memicu flag.}
    \begin{enumerate}
        \item Ubah \code{a\_lo} menjadi \hex{7FFF} dan \code{b\_lo}
              menjadi \hex{0001}.
        \item Rakit ulang, lalu periksa \flag{CF}, \flag{OF}, dan
              \flag{SF} setelah \asm{ADD}.
        \item Ulangi dengan \code{a\_lo} = \hex{8000} dan
              \code{b\_lo} = \hex{8000}.
        \item Jelaskan mengapa \flag{CF} dan \flag{OF} dapat berbeda.
    \end{enumerate}

    \item \textbf{Menguji divide overflow.}
    \begin{enumerate}
        \item Ubah pembagi pada bagian \asm{DIV} menjadi 3 dan
              nilai yang dibagi tetap 1000.
        \item Hitung hasil bagi seharusnya.
        \item Jelaskan mengapa program harus memeriksa \reg{AH}
              sebelum \asm{DIV}.
        \item Tunjukkan baris mana yang mencegah kesalahan tersebut.
    \end{enumerate}
\end{enumerate}

\subsection{Hasil yang Dikumpulkan}

\begin{enumerate}
    \item Berkas \file{operasi-aritmatika.asm} versi modifikasi.
    \item Tabel pengamatan register dan \eng{flag}.
    \item Tangkapan layar isi segmen data pada DEBUG.
    \item Analisis perbedaan \flag{CF} dan \flag{OF}.
    \item Penjelasan syarat aman pembagian.
\end{enumerate}

\begin{asesmenbox}
\vspace{1.1em}
\textbf{Asesmen praktikum 8 (bobot 5\% dari komponen praktikum).}

\begin{tabular}{@{}p{6.2cm}cp{4.6cm}@{}}
    \toprule
    \textbf{Indikator} & \textbf{Bobot} & \textbf{Kriteria} \\
    \midrule
    Ketepatan hasil ADD/ADC dan SUB/SBB & 30\% &
    Nilai 32 bit dan flag tercatat benar \\
    \addlinespace
    Ketepatan hasil MUL dan DIV & 25\% &
    Hasil bagi, sisa, dan letak register benar \\
    \addlinespace
    Analisis flag & 25\% &
    Perbedaan CF dan OF dijelaskan dengan tepat \\
    \addlinespace
    Analisis divide overflow & 20\% &
    Syarat aman dan baris pencegahannya dijelaskan \\
    \bottomrule
\end{tabular}
\end{asesmenbox}

\section{Latihan}
\label{sec:latihan-bab10}

\begin{enumerate}
    \item \textbf{ADD dan ADC.} Jelaskan perbedaan \asm{ADD} dan
          \asm{ADC}. Berikan satu contoh ketika keduanya menghasilkan
          nilai berbeda.

    \item \textbf{SUB dan SBB.} Jelaskan perbedaan \asm{SUB} dan
          \asm{SBB}. Mengapa \asm{SBB} diperlukan pada pengurangan
          bilangan lebih dari 16 bit?

    \item \textbf{INC dan DEC.} Mengapa \asm{INC} dan \asm{DEC}
          tidak boleh dipakai untuk meneruskan \eng{carry} pada
          aritmatika bertingkat?

    \item \textbf{Flag.} Tentukan \flag{CF}, \flag{OF}, \flag{ZF},
          dan \flag{SF} untuk setiap operasi 8 bit berikut.
          \begin{enumerate}
              \item \hex{7F} + \hex{01}
              \item \hex{80} + \hex{80}
              \item \hex{FF} + \hex{01}
          \end{enumerate}
          Jelaskan mengapa \flag{CF} dan \flag{OF} dapat berbeda.

    \item \textbf{Aritmatika 32 bit.} Diketahui dua bilangan 32 bit
          \hex{0001FFFF} dan \hex{00020001}. Tuliskan langkah
          penjumlahannya dengan \asm{ADD} dan \asm{ADC}, lalu
          tentukan hasil akhirnya.

    \item \textbf{MUL.} Jelaskan aturan operand \asm{MUL} untuk
          operand 8 bit dan 16 bit. Tentukan hasil dan letak hasil
          untuk:
          \begin{enumerate}
              \item \reg{AL} = 12, \reg{BL} = 5, \lstasm|MUL BL|
              \item \reg{AX} = 1000, \reg{BX} = 20, \lstasm|MUL BX|
          \end{enumerate}

    \item \textbf{DIV.} Jelaskan aturan operand \asm{DIV} untuk
          operand 8 bit dan 16 bit. Tentukan syarat aman agar tidak
          terjadi \eng{divide overflow}, lalu hitung hasil dan sisa
          dari $1000 \div 7$ dengan \lstasm|DIV BL|.

    \item \textbf{Penulisan program.} Tulislah program yang
          menghitung $(25 + 17) \times 4$, lalu menyimpan hasilnya
          pada variabel word. Gunakan \asm{ADD} dan \asm{MUL}, dan
          jelaskan letak hasil akhirnya.
\end{enumerate}

\begin{refleksibox}
\vspace{1.1em}
Jawab pertanyaan berikut pada jurnal praktikum Anda.
\begin{enumerate}
    \item Instruksi mana yang paling sulit Anda telusuri:
          \asm{ADC}/\asm{SBB}, \asm{MUL}, atau \asm{DIV}? Mengapa?
    \item Setelah memeriksa \eng{flag}, apakah Anda dapat membedakan
          luapan bertanda dan tidak bertanda? Jelaskan.
    \item Konsep apa yang perlu Anda ulangi sebelum mempelajari
          operasi logika dan \eng{shift} pada bab berikutnya?
\end{enumerate}
\end{refleksibox}

\section{Rangkuman}
\label{sec:rangkuman-bab10}

Pada bab ini telah dibahas:

\begin{itemize}
    \item \asm{ADD} menjumlahkan dua operand, sedangkan \asm{ADC}
          ikut menambahkan \flag{CF}. \asm{INC} menambah satu.

    \item \asm{SUB} mengurangi dua operand, sedangkan \asm{SBB}
          ikut mengurangkan \flag{CF}. \asm{DEC} mengurangi satu.

    \item Operasi aritmatika memperbarui \flag{CF}, \flag{OF},
          \flag{ZF}, \flag{SF}, \flag{AF}, dan \flag{PF}.

    \item \flag{CF} menandai luapan bilangan tidak bertanda,
          sedangkan \flag{OF} menandai luapan bilangan bertanda.

    \item \asm{INC} dan \asm{DEC} tidak mengubah \flag{CF}, sehingga
          tidak dipakai untuk meneruskan \eng{carry} pada
          aritmatika bertingkat.

    \item Bilangan lebih dari 16 bit dijumlahkan dengan \asm{ADD}
          pada word rendah dan \asm{ADC} pada word tinggi. Bentuk
          yang sama berlaku untuk pengurangan dengan \asm{SUB} dan
          \asm{SBB}.

    \item \asm{MUL} 8 bit mengalikan \reg{AL} dengan operand dan
          menyimpan hasil di \reg{AX}. \asm{MUL} 16 bit mengalikan
          \reg{AX} dengan operand dan menyimpan hasil di
          \reg{DX}:\reg{AX}.

    \item \asm{DIV} 8 bit membagi \reg{AX} dengan operand; hasil
          bagi di \reg{AL} dan sisa di \reg{AH}. \asm{DIV} 16 bit
          membagi \reg{DX}:\reg{AX}; hasil bagi di \reg{AX} dan
          sisa di \reg{DX}.

    \item \eng{Divide overflow} terjadi bila hasil bagi tidak muat
          pada register tujuan. Pemeriksaan \reg{AH} atau \reg{DX}
          harus dilakukan sebelum \asm{DIV}.
\end{itemize}

\begin{checklistbox}
\vspace{1.1em}
Tandai setiap butir yang sudah Anda kuasai.

\begin{itemize}
    \item[\checkbox] Saya dapat menggunakan \asm{ADD}, \asm{ADC},
          dan \asm{INC}.
    \item[\checkbox] Saya dapat menggunakan \asm{SUB}, \asm{SBB},
          dan \asm{DEC}.
    \item[\checkbox] Saya dapat menjelaskan pengaruh operasi
          aritmatika terhadap \eng{flag}.
    \item[\checkbox] Saya dapat membedakan \flag{CF} dan \flag{OF}.
    \item[\checkbox] Saya dapat menyusun aritmatika 32 bit dengan
          \asm{ADC} dan \asm{SBB}.
    \item[\checkbox] Saya dapat menggunakan \asm{MUL} sesuai aturan
          operand.
    \item[\checkbox] Saya dapat menggunakan \asm{DIV} dan memeriksa
          \eng{divide overflow}.
    \item[\checkbox] Saya dapat memverifikasi hasil dengan DEBUG
          atau emu8086.
\end{itemize}

\vspace{0.4em}
Bila masih ada butir yang belum bertanda, ulangi bagian yang
bersesuaian sebelum melanjutkan ke bab berikutnya.
\end{checklistbox}

\begin{infobox}
\vspace{1.1em}
\textbf{Rujukan bab.} Materi pada bab ini diringkas dari
\cite{gunawan2019}, Bagian 14, hlm.~127--142, dan
\cite{sto2001}, BAB VIII, hlm.~32--40.
\end{infobox}

\section{Evaluasi}
\label{sec:evaluasi-bab10}

\subsection{Soal Pilihan Ganda}

\begin{enumerate}
    \item Perbedaan utama \asm{ADC} dari \asm{ADD} adalah \dots
          \begin{enumerate}
              \item \asm{ADC} hanya bekerja pada 8 bit
              \item \asm{ADC} ikut menambahkan \flag{CF}
              \item \asm{ADC} tidak mengubah \eng{flag}
              \item \asm{ADC} menyimpan hasil di memori
          \end{enumerate}

    \item Instruksi yang tidak mengubah \flag{CF} adalah \dots
          \begin{enumerate}
              \item \asm{ADD} \item \asm{SUB}
              \item \asm{INC} \item \asm{ADC}
          \end{enumerate}

    \item Operasi \hex{7F} + \hex{01} pada 8 bit menghasilkan
          \dots
          \begin{enumerate}
              \item \flag{CF} = 1, \flag{OF} = 0
              \item \flag{CF} = 0, \flag{OF} = 1
              \item \flag{CF} = 1, \flag{OF} = 1
              \item \flag{CF} = 0, \flag{OF} = 0
          \end{enumerate}

    \item Hasil \asm{MUL} 8 bit diletakkan pada \dots
          \begin{enumerate}
              \item \reg{AL} \item \reg{AH}
              \item \reg{AX} \item \reg{DX}:\reg{AX}
          \end{enumerate}

    \item Sisa hasil \asm{DIV} 8 bit diletakkan pada \dots
          \begin{enumerate}
              \item \reg{AL} \item \reg{AH}
              \item \reg{AX} \item \reg{DX}
          \end{enumerate}
\end{enumerate}

\subsection{Soal Uraian}

\begin{enumerate}
    \item Jelaskan perbedaan \asm{ADD} dan \asm{ADC}, serta
          \asm{SUB} dan \asm{SBB}. Tuliskan contoh penjumlahan
          32 bit yang memerlukan \asm{ADC}.

    \item Jelaskan perbedaan \flag{CF} dan \flag{OF}. Berikan
          contoh 8 bit yang menunjukkan \flag{CF} = 1 tetapi
          \flag{OF} = 0, dan contoh \flag{OF} = 1 tetapi
          \flag{CF} = 0.

    \item Jelaskan aturan operand \asm{MUL} dan \asm{DIV} untuk
          ukuran 8 bit dan 16 bit. Tuliskan pula syarat aman agar
          tidak terjadi \eng{divide overflow}.

    \item Tulislah program yang menghitung $(25 + 17) \times 4$,
          lalu menyimpan hasilnya pada variabel word. Jelaskan
          register yang dipakai dan letak hasil akhirnya.
\end{enumerate}

\subsection{Kisi-kisi Penilaian}

\begin{table}[H]
    \centering
    \caption{Kisi-kisi evaluasi Bab~\ref{chap:aritmatika}.}
    \label{tab:kisi-bab10}
    \small
    \begin{tabular}{@{}p{1.5cm}p{5.0cm}p{3.0cm}c@{}}
        \toprule
        \textbf{Sub-CPMK} & \textbf{Indikator yang diukur} &
        \textbf{Bentuk soal} & \textbf{Bobot} \\
        \midrule
        3.1 & Menerapkan ADD, ADC, INC, SUB, SBB, DEC &
        Pilihan ganda 1, 2; Uraian 1 & 30\% \\
        \addlinespace
        3.1 & Menjelaskan pengaruh operasi terhadap flag &
        Pilihan ganda 3; Uraian 2 & 25\% \\
        \addlinespace
        3.1 & Menerapkan MUL dan DIV serta overflow &
        Pilihan ganda 4, 5; Uraian 3 & 30\% \\
        \addlinespace
        3.1 & Menulis program aritmatika &
        Uraian 4 & 15\% \\
        \bottomrule
    \end{tabular}
\end{table}

\begin{checklistbox}
\vspace{1.1em}
\textbf{Checklist Pencapaian Kompetensi
Bab~\ref{chap:aritmatika}}

Seorang mahasiswa dinyatakan kompeten pada Sub-CPMK 3.1 apabila:
\begin{itemize}
    \item[\checked] dapat menggunakan ADD, ADC, dan INC;
    \item[\checked] dapat menggunakan SUB, SBB, dan DEC;
    \item[\checked] dapat menjelaskan pengaruh operasi terhadap
          \eng{flag};
    \item[\checked] dapat membedakan CF dan OF;
    \item[\checked] dapat menyusun aritmatika 32 bit;
    \item[\checked] dapat menggunakan MUL sesuai aturan operand;
    \item[\checked] dapat menggunakan DIV dan memeriksa overflow;
    \item[\checked] dapat memverifikasi hasil dengan debugger.
\end{itemize}
\end{checklistbox}


\end{document}
